\begin{figure}[H]
	\centering 
	\includegraphics[width=1.0\textwidth]{SpecificaTecnica/Chat/ChatBackend.png}
	\caption{Diagramma delle classi per la gestione del carrello}
\end{figure}

\paragraph{Descrizione generale}
Questa parte si occupa della gestione del carrello associato a uno specifico utente.\\
Le classi descritte consentono di recuperare i prodotti presenti nel carrello, aggiungere nuovi articoli, rimuoverli, modificarne la quantit\`a e inviare l'ordine finale.\\
L'architettura segue il pattern Ports and Adapters, separando la logica applicativa dai dettagli di persistenza e integrando un servizio dedicato alla ricerca semantica dei prodotti presenti nel carrello.

\subsubsubsection{EmbeddedCartService}
\label{EmbeddedCartService}
\begin{figure}[H]
	\centering 
	\includegraphics[width=0.9\textwidth]{SpecificaTecnica/Chat/EmbeddedCartService.png}
	\caption{Gestione del carrello: EmbeddedCartService}
\end{figure}

\paragraph{Descrizione}
\textit{EmbeddedCartService} \`e il service applicativo responsabile del caricamento del contenuto del carrello in forma vettoriale e della successiva ricerca semantica.\\
Recupera i prodotti del carrello tramite la porta \textit{CartRepoPort}, ne costruisce una rappresentazione vettoriale tramite \textit{EmbedderPort} e interroga il database vettoriale attraverso \textit{VecDbPortOut}.

\paragraph{Attributi}
\begin{itemize}
	\item \textbf{cart\_vect: VecDbPortOut} - attributo privato di tipo \textit{VecDbPortOut}.\\
	Porta utilizzata per memorizzare e interrogare la rappresentazione vettoriale dei prodotti presenti nel carrello.
	\item \textbf{cart\_repo: CartRepoPort} - attributo privato di tipo \textit{CartRepoPort}.\\
	Porta utilizzata per recuperare i prodotti presenti nel carrello dell'utente.
	\item \textbf{embedder: EmbedderPort} - attributo privato di tipo \textit{EmbedderPort}.\\
	Porta utilizzata per generare gli embedding dei prodotti da indicizzare nel database vettoriale.
\end{itemize}

\paragraph{Metodi e funzioni}
\begin{itemize}
	\item \textbf{\_\_init\_\_(cart\_vect: VecDbPortOut, cart\_repo: CartRepoPort, embedder: EmbedderPort)} - metodo pubblico di inizializzazione.\\
	Riceve e memorizza le dipendenze necessarie per il recupero dei dati del carrello, la generazione degli embedding e l'interrogazione vettoriale.
	\item \textbf{load\_cart(username: string)} - metodo pubblico che recupera dal repository i prodotti presenti nel carrello dell'utente identificato da \textit{username} e li carica nel sistema di ricerca vettoriale.\\
	Questo passaggio prepara i dati necessari per successive ricerche semantiche.
	\item \textbf{search\_cart(username: string, query: string, threshold: float): string*} - metodo pubblico che esegue una ricerca semantica all'interno del carrello dell'utente, utilizzando la query fornita e una soglia di similarit\`a.\\
	Restituisce gli identificativi o i riferimenti ai prodotti che soddisfano i criteri di ricerca.
\end{itemize}

\subsubsubsection{ToolCartService}
\begin{figure}[H]
	\centering 
	\includegraphics[width=1.0\textwidth]{SpecificaTecnica/Chat/ToolCartService.png}
	\caption{Gestione del carrello: ToolCartService}
\end{figure}

\paragraph{Descrizione}
\textit{ToolCartService} \`e il service applicativo che espone le principali operazioni di business sul carrello dell'utente.\\
Si occupa di coordinare l'accesso ai dati tramite \textit{CartRepoPort} e di sfruttare \textit{EmbeddedCartService} per le ricerche semantiche sui prodotti gi\`a presenti.

\paragraph{Attributi}
\begin{itemize}
	\item \textbf{username: string} - attributo privato di tipo \textit{string}.\\
	Rappresenta lo username dell'utente a cui il carrello \`e associato.
	\item \textbf{cart\_repo: CartRepoPort} - attributo privato di tipo \textit{CartRepoPort}.\\
	Porta utilizzata per eseguire le operazioni di lettura e aggiornamento del carrello.
	\item \textbf{embedded\_cart: EmbeddedCartService} - attributo privato di tipo \textit{EmbeddedCartService}.\\
	Service utilizzato per effettuare ricerche semantiche tra i prodotti gi\`a presenti nel carrello.
\end{itemize}

\paragraph{Metodi e funzioni}
\begin{itemize}
	\item \textbf{\_\_init\_\_(username: string, cart\_repo: CartRepoPort, embedded\_cart: EmbeddedCartService)} - metodo pubblico di inizializzazione.\\
	Riceve lo username dell'utente e le dipendenze necessarie alla gestione del carrello.
	\item \textbf{get\_cart\_items(): CartProducts*} - metodo pubblico che recupera l'insieme dei prodotti attualmente presenti nel carrello dell'utente.\\
	Restituisce una collezione di oggetti di tipo \textit{CartProduct}.
	\item \textbf{search\_cart(query: string, threshold: float): CartProduct*} - metodo pubblico che ricerca prodotti nel carrello in base a una query semantica.\\
	Delega la ricerca a \textit{EmbeddedCartService} e restituisce i prodotti corrispondenti.
	\item \textbf{add\_to\_cart(prod\_id: string, qty: int): CartProduct} - metodo pubblico che aggiunge al carrello il prodotto identificato da \textit{prod\_id} nella quantit\`a indicata.\\
	Restituisce il prodotto aggiornato nel carrello.
	\item \textbf{remove\_from\_cart(prod\_id: string): CartProduct} - metodo pubblico che rimuove dal carrello il prodotto identificato da \textit{prod\_id}.\\
	Restituisce il prodotto rimosso o il relativo stato aggiornato.
	\item \textbf{update\_cart\_item\_qty(prod\_id: string, qty: int, operation: CartUpdateOperation): CartProduct} - metodo pubblico che aggiorna la quantit\`a di un prodotto gi\`a presente nel carrello.\\
	L'aggiornamento viene effettuato in base all'operazione specificata da \textit{CartUpdateOperation}.
\end{itemize}

\subsubsubsection{CartService}
\begin{figure}[H]
	\centering 
	\includegraphics[width=0.7\textwidth]{SpecificaTecnica/Chat/CartService.png}
	\caption{Gestione del carrello: CartService}
\end{figure}

\paragraph{Descrizione}
\textit{CartService} \`e il service applicativo dedicato alle operazioni tradizionali di gestione del carrello e all'invio dell'ordine.\\
Rappresenta un punto di accesso semplificato alle funzionalit\`a di lettura e modifica del carrello, delegando le operazioni alla porta \textit{CartRepoPort}.

\paragraph{Attributi}
\begin{itemize}
	\item \textbf{adapter: CartRepoPort} - attributo privato di tipo \textit{CartRepoPort}.\\
	Porta utilizzata per accedere alle operazioni di persistenza del carrello.
\end{itemize}

\paragraph{Metodi e funzioni}
\begin{itemize}
	\item \textbf{\_\_init\_\_(adapter: CartRepoPort)} - metodo pubblico di inizializzazione.\\
	Riceve e memorizza la porta di accesso al repository del carrello.
	\item \textbf{get\_cart\_products(username: string): CartProduct*} - metodo pubblico che recupera tutti i prodotti presenti nel carrello dell'utente identificato da \textit{username}.\\
	Restituisce una collezione di oggetti di tipo \textit{CartProduct}.
	\item \textbf{send\_order(username: string)} - metodo pubblico che avvia l'invio dell'ordine relativo al carrello dell'utente specificato.\\
	Delega la persistenza finale dell'operazione alla porta \textit{CartRepoPort}.
\end{itemize}

\subsubsubsection{CartRepoPort}
\begin{figure}[H]
	\centering 
	\includegraphics[width=1.0\textwidth]{SpecificaTecnica/Chat/CartRepoPort.png}
	\caption{Gestione del carrello: CartRepoPort}
\end{figure}

\paragraph{Descrizione}
\textit{CartRepoPort} \`e un'interfaccia che definisce il contratto per l'accesso e la modifica dei dati relativi al carrello.\\
Viene implementata da \textit{CartRepoAdapter} e utilizzata dai service applicativi per mantenere il disaccoppiamento tra logica di business e persistenza.

\paragraph{Attributi}
\textit{CartRepoPort} non definisce attributi.

\paragraph{Metodi e funzioni}
\begin{itemize}
	\item \textbf{get\_products(username: string): CartProduct*} - metodo pubblico che recupera i prodotti presenti nel carrello dell'utente identificato da \textit{username}.\\
	Restituisce una collezione di oggetti di tipo \textit{CartProduct}.
	\item \textbf{add\_product(prod\_id: string, username: string, qty: int): CartProduct} - metodo pubblico che aggiunge al carrello dell'utente il prodotto identificato da \textit{prod\_id} nella quantit\`a indicata.\\
	Restituisce il prodotto aggiunto o aggiornato.
	\item \textbf{remove\_product(prod\_id: string, username: string): CartProduct} - metodo pubblico che rimuove dal carrello dell'utente il prodotto identificato da \textit{prod\_id}.\\
	Restituisce il prodotto coinvolto nell'operazione.
	\item \textbf{update\_quantity(prod\_id: string, username: string, qty: int, operation: CartUpdateOperation): CartProduct} - metodo pubblico che aggiorna la quantit\`a di un prodotto gi\`a presente nel carrello.\\
	L'aggiornamento avviene secondo l'operazione specificata da \textit{CartUpdateOperation}.
	\item \textbf{send\_order(username: string): void} - metodo pubblico che conferma e inoltra l'ordine associato al carrello dell'utente.
\end{itemize}

\subsubsubsection{CartRepoAdapter}
\begin{figure}[H]
	\centering 
	\includegraphics[width=0.9\textwidth]{SpecificaTecnica/Chat/CartRepoAdapter.png}
	\caption{Gestione del carrello: CartRepoAdapter}
\end{figure}

\paragraph{Descrizione}
\textit{CartRepoAdapter} \`e la classe concreta che implementa l'interfaccia \textit{CartRepoPort}.\\
Agisce da adattatore tra i service applicativi e il repository concreto \textit{CartRepository}, traducendo le operazioni richieste dal dominio in chiamate al layer di persistenza.

\paragraph{Attributi}
\begin{itemize}
	\item \textbf{repo: CartRepository} - attributo privato di tipo \textit{CartRepository}.\\
	Rappresenta il repository concreto utilizzato per eseguire le operazioni sul database.
\end{itemize}

\paragraph{Metodi e funzioni}
Oltre ad essere implementate tutte le funzioni di \textit{CartRepoPort}, \`e disponibile anche la seguente funzione:
\begin{itemize}
	\item \textbf{\_\_init\_\_(repo: CartRepository)} - metodo pubblico di inizializzazione.\\
	Riceve e memorizza il repository concreto utilizzato per l'accesso ai dati del carrello.
\end{itemize}

\subsubsubsection{CartRepository}
\begin{figure}[H]
	\centering 
	\includegraphics[width=1.0\textwidth]{SpecificaTecnica/Chat/CartRepository.png}
	\caption{Gestione del carrello: CartRepository}
\end{figure}

\paragraph{Descrizione}
\textit{CartRepository} \`e la classe che gestisce direttamente le operazioni di lettura e scrittura dei dati del carrello sul database.\\
Utilizza una sessione di database per eseguire il recupero dei prodotti, l'aggiunta di nuovi articoli, la rimozione di elementi, l'aggiornamento delle quantit\`a e l'invio finale dell'ordine.

\paragraph{Attributi}
\begin{itemize}
	\item \textbf{db: Session} - attributo privato di tipo \textit{Session}.\\
	Rappresenta la sessione di database utilizzata per eseguire le operazioni di persistenza.
\end{itemize}

\paragraph{Metodi e funzioni}
\begin{itemize}
	\item \textbf{\_\_init\_\_(db: Session)} - metodo pubblico di inizializzazione.\\
	Riceve una sessione di database e la memorizza per l'esecuzione delle operazioni sul carrello.
	\item \textbf{get\_products(username: string): CartProductRepository*} - metodo pubblico che recupera dal database i prodotti presenti nel carrello dell'utente indicato.\\
	Restituisce una collezione di entit\`a repository associate ai prodotti del carrello.
	\item \textbf{add\_product(prod\_id: string, username: string, qty: int): CartProductRepository} - metodo pubblico che inserisce nel database un nuovo prodotto nel carrello dell'utente oppure ne aggiorna la presenza secondo la logica prevista.\\
	Restituisce l'entit\`a repository risultante.
	\item \textbf{remove\_product(username: string, prod\_id: string): CartProductRepository} - metodo pubblico che rimuove dal database il prodotto indicato dal carrello dell'utente.\\
	Restituisce l'entit\`a repository coinvolta nell'operazione.
	\item \textbf{update\_quantity(username: string, prod\_id: string, qty: int, operation: CartUpdateOperation): CartProductRepository} - metodo pubblico che aggiorna nel database la quantit\`a del prodotto indicato nel carrello dell'utente.\\
	La modifica viene applicata in accordo all'operazione specificata.
	\item \textbf{send\_order(username: string)} - metodo pubblico che conclude il processo di acquisto inviando l'ordine associato al carrello dell'utente.
\end{itemize}

\subsubsubsection{Boh, parte separata non so come volete dividere}
\begin{figure}[H]
	\centering 
	\includegraphics[width=1.0\textwidth]{SpecificaTecnica/Chat/UMLChatService.png}
	\caption{Diagramma delle classi per la gestione della chat}
\end{figure}

\paragraph{Descrizione generale}
Questa parte si occupa della gestione delle conversazioni tra utente e sistema.\\
Le classi descritte consentono di recuperare i messaggi di una conversazione, accedere allo storico delle chat, creare nuove conversazioni e aggiungere nuovi messaggi.\\
Anche in questo caso l'architettura segue il pattern Ports and Adapters, mantenendo separata la logica applicativa dalla persistenza dei dati.

\subsubsubsection{ChatService}
\label{ChatService}
\begin{figure}[H]
	\centering 
	\includegraphics[width=0.9\textwidth]{SpecificaTecnica/Chat/ChatService.png}
	\caption{Gestione della chat: ChatService}
\end{figure}

\paragraph{Descrizione}
\textit{ChatService} \`e il service applicativo responsabile della gestione delle conversazioni e dei messaggi scambiati nella chat.\\
Coordina l'accesso ai dati tramite \textit{ChatRepoPort} e integra eventuali dipendenze aggiuntive, come il componente LLM e il servizio di storico, per supportare il flusso conversazionale.

\paragraph{Attributi}
\begin{itemize}
	\item \textbf{repo: ChatRepoPort} - attributo privato di tipo \textit{ChatRepoPort}.\\
	Porta utilizzata per accedere ai dati delle conversazioni e dei messaggi.
	\item \textbf{llm: LLMPort} - attributo privato di tipo \textit{LLMPort}.\\
	Porta utilizzata per interagire con il modello linguistico durante la gestione della conversazione.
	\item \textbf{storico\_service: StoricoService | Null} - attributo privato che rappresenta un eventuale riferimento al servizio di gestione dello storico ordini.\\
	Viene utilizzato nei casi in cui la conversazione debba integrarsi con funzionalit\`a legate allo storico.
\end{itemize}

\paragraph{Metodi e funzioni}
\begin{itemize}
	\item \textbf{\_\_init\_\_(repo: ChatRepoPort, llm: LLMPort, storico\_service: StoricoService | Null)} - metodo pubblico di inizializzazione.\\
	Riceve e memorizza le dipendenze necessarie alla gestione della chat.
	\item \textbf{get\_all\_messages(conv\_id: int): Message*} - metodo pubblico che recupera tutti i messaggi associati alla conversazione identificata da \textit{conv\_id}.\\
	Restituisce una collezione di oggetti di tipo \textit{Message}.
	\item \textbf{send\_message(conv\_id: int, username: string, content: string): Message} - metodo pubblico che gestisce l'invio di un nuovo messaggio all'interno della conversazione specificata.\\
	Il metodo coordina il salvataggio del messaggio e la produzione della risposta associata al flusso conversazionale.
\end{itemize}

\subsubsubsection{ChatRepoPort}
\begin{figure}[H]
	\centering 
	\includegraphics[width=1.0\textwidth]{SpecificaTecnica/Chat/ChatRepoPort.png}
	\caption{Gestione della chat: ChatRepoPort}
\end{figure}

\paragraph{Descrizione}
\textit{ChatRepoPort} \`e un'interfaccia che definisce il contratto per l'accesso ai dati relativi alle conversazioni e ai messaggi della chat.\\
Viene implementata da \textit{ChatRepoAdapter} e utilizzata da \textit{ChatService} per garantire il disaccoppiamento tra logica applicativa e persistenza.

\paragraph{Attributi}
\textit{ChatRepoPort} non definisce attributi.

\paragraph{Metodi e funzioni}
\begin{itemize}
	\item \textbf{get\_messages(conv\_id: int): Message*} - metodo pubblico che recupera i messaggi della conversazione identificata da \textit{conv\_id}.\\
	Restituisce una collezione di oggetti di tipo \textit{Message}.
	\item \textbf{get\_chat\_history(conv\_id: int, max\_messages: int): Message*} - metodo pubblico che recupera una porzione dello storico della conversazione, limitata al numero massimo di messaggi indicato.\\
	Restituisce una collezione di oggetti di tipo \textit{Message}.
	\item \textbf{add\_message(conv\_id: int, text: string, sender: SenderEnum): Message} - metodo pubblico che aggiunge un nuovo messaggio alla conversazione specificata.\\
	Restituisce il messaggio appena salvato.
	\item \textbf{create\_conversation(username: string): Conversation} - metodo pubblico che crea una nuova conversazione associata all'utente indicato.\\
	Restituisce l'oggetto \textit{Conversation} creato.
	\item \textbf{conversation\_exist(conv\_id: int): bool} - metodo pubblico che verifica l'esistenza di una conversazione identificata da \textit{conv\_id}.\\
	Restituisce \textit{true} se la conversazione esiste, \textit{false} altrimenti.
\end{itemize}

\subsubsubsection{ChatRepoAdapter}
\begin{figure}[H]
	\centering 
	\includegraphics[width=0.9\textwidth]{SpecificaTecnica/Chat/ChatRepoAdapter.png}
	\caption{Gestione della chat: ChatRepoAdapter}
\end{figure}

\paragraph{Descrizione}
\textit{ChatRepoAdapter} \`e la classe concreta che implementa l'interfaccia \textit{ChatRepoPort}.\\
Funge da adattatore tra il livello applicativo e il repository concreto \textit{ChatRepository}, delegando a quest'ultimo l'esecuzione delle operazioni di persistenza.

\paragraph{Attributi}
\begin{itemize}
	\item \textbf{repo: ChatRepository} - attributo privato di tipo \textit{ChatRepository}.\\
	Rappresenta il repository concreto utilizzato per la gestione dei dati della chat.
\end{itemize}

\paragraph{Metodi e funzioni}
Oltre ad essere implementate tutte le funzioni di \textit{ChatRepoPort}, \`e disponibile anche la seguente funzione:
\begin{itemize}
	\item \textbf{\_\_init\_\_(repo: ChatRepository)} - metodo pubblico di inizializzazione.\\
	Riceve e memorizza il repository concreto utilizzato per accedere ai dati delle conversazioni.
\end{itemize}

\subsubsubsection{ChatRepository}
\begin{figure}[H]
	\centering 
	\includegraphics[width=1.0\textwidth]{SpecificaTecnica/Chat/ChatRepository.png}
	\caption{Gestione della chat: ChatRepository}
\end{figure}

\paragraph{Descrizione}
\textit{ChatRepository} \`e la classe che gestisce direttamente le operazioni di lettura e scrittura delle conversazioni e dei messaggi sul database.\\
Utilizza una sessione di database per creare nuove conversazioni, recuperare i messaggi associati e salvare i nuovi contenuti prodotti durante l'interazione in chat.

\paragraph{Attributi}
\begin{itemize}
	\item \textbf{db: Session} - attributo privato di tipo \textit{Session}.\\
	Rappresenta la sessione di database utilizzata per eseguire le operazioni di persistenza della chat.
\end{itemize}

\paragraph{Metodi e funzioni}
\begin{itemize}
	\item \textbf{\_\_init\_\_(db: Session)} - metodo pubblico di inizializzazione.\\
	Riceve una sessione di database e la memorizza per l'esecuzione delle operazioni sulle conversazioni.
	\item \textbf{get\_conversation(conv\_id: int): Conversazioni | Null} - metodo pubblico che recupera la conversazione identificata da \textit{conv\_id}.\\
	Restituisce la conversazione trovata oppure \textit{Null} nel caso in cui non esista.
	\item \textbf{create\_conversation(username: string): Conversation} - metodo pubblico che crea e persiste una nuova conversazione associata allo username fornito.\\
	Restituisce l'entit\`a relativa alla conversazione creata.
	\item \textbf{get\_messages(conv\_id: int): Message*} - metodo pubblico che recupera tutti i messaggi della conversazione indicata.\\
	Restituisce una collezione di messaggi.
	\item \textbf{get\_chat\_history(conv\_id: int, text: string, sender: SenderEnum): Message*} - metodo pubblico che recupera lo storico della chat secondo i criteri previsti dall'implementazione.\\
	Restituisce una collezione di messaggi utile alla ricostruzione del contesto conversazionale.
	\item \textbf{add\_message(conv\_id: int, text: string, sender: SenderEnum): Message} - metodo pubblico che salva un nuovo messaggio nella conversazione indicata.\\
	Restituisce il messaggio persistito.
\end{itemize}
